\documentclass[
    language=english,
    doctype=technical-report,
    institution=none,
    titlestyle=book,
    oneside,
]{../../omnilatex}

\RequirePackage{config/document-settings}

\addbibresource{../../bib/bibliography.bib}

\begin{document}

\maketitle

\tableofcontents

%=============================================================================
\chapter{Introduction}
\label{sec:introduction}
%=============================================================================

This document describes the system architecture of DataPipe, a streaming
ETL platform. It is intended for architects, senior engineers, and
technical stakeholders who need to understand the system's design
rationale, component interactions, and technology choices.

\section{Document Scope}

This architecture document covers:

\begin{itemize}
    \item High-level system architecture and component boundaries.
    \item Data flow through the platform.
    \item Technology selection rationale (Architecture Decision Records).
    \item Non-functional requirements and their architectural implications.
    \item Integration patterns with external systems.
\end{itemize}

\section{Architecture Principles}

\begin{enumerate}
    \item \textbf{Separation of concerns:} Each component has a single,
        well-defined responsibility.
    \item \textbf{Extensibility:} New connectors and transforms can be
        added without modifying the core runtime.
    \item \textbf{Resilience:} The system tolerates component failures
        without data loss.
    \item \textbf{Observability:} All components expose metrics, logs,
        and traces.
    \item \textbf{Simplicity:} Prefer boring technology over novel
        solutions.
\end{enumerate}

%=============================================================================
\chapter{System Architecture}
\label{sec:architecture}
%=============================================================================

\section{Component Overview}

The system comprises six major components arranged in a layered
architecture:

\begin{figure}[htbp]
    \centering
    \begin{tikzpicture}[
        node distance=1cm and 1.5cm,
        layer/.style={draw, rounded corners, minimum width=10cm,
                      minimum height=1.5cm, text centered,
                      fill=#1!10, draw=#1!50!black,
                      font=\small},
        comp/.style={draw, rounded corners, minimum width=2.2cm,
                     minimum height=0.7cm, text centered, font=\footnotesize,
                     fill=white, thick},
        arr/.style={->, >=stealth, thick},
        annot/.style={font=\scriptsize, text width=2.5cm, align=center},
    ]
        \node[layer=blue] (api) {API Layer};
        \node[comp] (rest) [above=0.3cm of api] {REST API};
        \node[comp] (cli) [left=1cm of rest] {CLI};
        \node[comp] (grpc) [right=1cm of rest] {gRPC};

        \node[layer=green] (runtime) [below=0.6cm of api] {Core Runtime};
        \node[comp] (pipeline) [above=0.3cm of runtime] {Pipeline Engine};
        \node[comp] (scheduler) [left=0.5cm of pipeline] {Scheduler};
        \node[comp] (executor) [right=0.5cm of pipeline] {Executor};

        \node[layer=orange] (connectors) [below=0.6cm of runtime] {Connector Layer};
        \node[comp] (source) [above=0.3cm of connectors] {Sources};
        \node[comp] (sink) [right=2cm of source] {Sinks};
        \node[comp] (window) [left=2cm of source] {Windows};

        \node[layer=red] (storage) [below=0.6cm of connectors] {Storage Layer};
        \node[comp] (kafka) [above=0.3cm of storage] {Kafka};
        \node[comp] (pg) [right=2cm of kafka] {PostgreSQL};
        \node[comp] (s3) [left=2cm of kafka] {S3};

        \draw[arr] (api) -- (runtime);
        \draw[arr] (runtime) -- (connectors);
        \draw[arr] (connectors) -- (storage);
    \end{tikzpicture}
    \caption{DataPipe layered architecture.}
    \label{fig:layered-arch}
\end{figure}

\section{Data Flow}

The following diagram shows how a record flows through the system from
source to sink:

\begin{figure}[htbp]
    \centering
    \begin{tikzpicture}[
        node distance=1.5cm and 2cm,
        box/.style={draw, rounded corners, minimum width=2.5cm,
                     minimum height=0.9cm, text centered, font=\small,
                     fill=#1!20, draw=#1!60!black},
        arr/.style={->, >=stealth, thick},
    ]
        \node[box=blue]   (src)   {Source};
        \node[box=green]  ( deser) [right=of src] {Deserialize};
        \node[box=orange] (trans) [right=of deser] {Transform};
        \node[box=purple] (ser)   [right=of trans] {Serialize};
        \node[box=red]    (sink)  [right=of ser] {Sink};

        \draw[arr] (src) -- (deser);
        \draw[arr] (deser) -- (trans);
        \draw[arr] (trans) -- (ser);
        \draw[arr] (ser) -- (sink);

        \node[font=\footnotesize, below=0.3cm of deser, text width=2cm, align=center]
            {JSON/Avro/Protobuf};
        \node[font=\footnotesize, below=0.3cm of trans, text width=2cm, align=center]
            {Filter/Map/Aggregate};
        \node[font=\footnotesize, below=0.3cm of ser, text width=2cm, align=center]
            {Target format};
    \end{tikzpicture}
    \caption{Record processing flow.}
    \label{fig:data-flow}
\end{figure}

%=============================================================================
\chapter{Technology Decisions}
\label{sec:technology}
%=============================================================================

\section{Language and Runtime}

\begin{itemize}
    \item \textbf{Java 21} chosen for the core runtime due to its mature
        ecosystem, strong typing, and excellent Kafka/Flink integration.
    \item \textbf{Python 3.12} chosen for the CLI and custom transforms
        due to rapid development cycle and data science community.
\end{itemize}

\section{Messaging}

\textbf{Apache Kafka} was selected over RabbitMQ and NATS for the
following reasons:

\begin{enumerate}
    \item Persistent storage enables replay and exactly-once semantics.
    \item Partition-based parallelism scales horizontally.
    \item Mature ecosystem with connectors for most data sources.
    \item Strong community and commercial support.
\end{enumerate}

\section{Database}

\textbf{PostgreSQL 16} serves as the primary metadata store:

\begin{itemize}
    \item ACID compliance ensures pipeline configuration consistency.
    \item JSONB support enables flexible metadata storage.
    \item Replication provides high availability.
    \item Rich extension ecosystem (pg\_stat\_statements, pg\_cron).
\end{itemize}

%=============================================================================
\chapter{Architecture Decision Records}
\label{sec:adr}
%=============================================================================

\section{ADR-001: Use Kafka as the Primary Message Broker}

\begin{description}
    \item[Status] Accepted.
    \item[Context] The platform requires a durable, high-throughput message
        broker for inter-component communication and data buffering.
    \item[Decision] Use Apache Kafka as the primary message broker.
    \item[Consequences] Positive: durable storage, replay capability,
        exactly-once semantics. Negative: operational complexity, JVM
        dependency.
\end{description}

\section{ADR-002: Plugin Architecture for Connectors}

\begin{description}
    \item[Status] Accepted.
    \item[Context] New data sources and sinks must be added without
        modifying the core runtime.
    \item[Decision] Implement a Java ServiceLoader-based plugin system for
        connectors.
    \item[Consequences] Positive: clean separation, third-party
        extensibility. Negative: additional abstraction layer, versioning
        complexity.
\end{description}

\section{ADR-003: Python CLI Over Go or Rust}

\begin{description}
    \item[Status] Accepted.
    \item[Context] The CLI needs rapid iteration and should be accessible
        to data engineers who primarily use Python.
    \item[Decision] Implement the CLI in Python using Click.
    \item[Consequences] Positive: fast development, familiar to target
        users. Negative: slower startup, deployment complexity.
\end{description}

%=============================================================================
\chapter{Non-Functional Requirements}
\label{sec:nfr}
%=============================================================================

\begin{table}[htbp]
    \centering
    \caption{Non-functional requirements and architectural implications.}
    \label{tab:nfr}
    \begin{tabular}{@{}p{3.5cm}p{4cm}p{5cm}@{}}
        \toprule
        Requirement & Target & Architectural Implication \\
        \midrule
        Throughput & 500k events/s & Horizontal scaling of pipeline
            executor nodes; Kafka partitioning. \\
        Latency (P99) & $< 200$\,ms & In-memory transforms; avoid
            synchronous I/O in the hot path. \\
        Availability & 99.95\% & Multi-AZ deployment; automated failover;
            read replicas. \\
        Durability & At-least-once & Kafka WAL; idempotent sink writes;
            checkpointing. \\
        Scalability & 10x growth & Stateless executors; partition-based
            parallelism; auto-scaling. \\
        \bottomrule
    \end{tabular}
\end{table}

%=============================================================================
\chapter{Integration Architecture}
\label{sec:integration}
%=============================================================================

\section{External System Integrations}

DataPipe integrates with the following external systems:

\begin{figure}[htbp]
    \centering
    \begin{tikzpicture}[
        node distance=1.2cm and 1.5cm,
        box/.style={draw, rounded corners, minimum width=2.4cm,
                     minimum height=0.8cm, text centered, font=\small},
        center/.style={box, fill=red!20, draw=red!60!black,
                       minimum width=3cm},
        arr/.style={->, >=stealth, thick},
    ]
        \node[box=blue]   (pg)     {PostgreSQL};
        \node[box=green]  (s3)     [below=of pg] {S3};
        \node[box=orange] (bigq)   [below=of s3] {BigQuery};
        \node[center]     (dp)     [right=2cm of s3] {DataPipe};
        \node[box=purple] (kafka)  [right=2cm of pg] {Kafka};
        \node[box=cyan]   (redis)  [below=of kafka] {Redis};
        \node[box=red]    (elastic)[below=of redis] {Elasticsearch};

        \draw[arr] (pg) -- (dp);
        \draw[arr] (s3) -- (dp);
        \draw[arr] (bigq) -- (dp);
        \draw[arr] (dp) -- (kafka);
        \draw[arr] (dp) -- (redis);
        \draw[arr] (dp) -- (elastic);
    \end{tikzpicture}
    \caption{External system integration map.}
    \label{fig:integration}
\end{figure}

\section{API Versioning Strategy}

The REST API follows semantic versioning:

\begin{lstlisting}[language=bash, caption={API versioning examples.}]
# Versioned endpoints
/api/v1/pipelines
/api/v2/pipelines

# Content negotiation via Accept header
curl -H "Accept: application/vnd.datapipe.v2+json" \\
    http://localhost:8080/api/v2/pipelines
\end{lstlisting}

\printbibliography[heading=bibintoc, title={References}]

\end{document}
